\chapter{皮むき動作の実現}
\label{chap:peeling}

% 第一章
\section{皮むき動作の概要}
\label{sec: peeling-overview}

% 第一段写了几句废话啊。。。
本章は，第1章で提示した3Dモデルを基盤とした分析主導型アプローチを，皮むき動作における幾何的な軌道生成と力覚に基づく接触維持へ具体化するものである．第1章の表1.2では，皮むき動作を，法線・軌道に関する幾何学的分析と，実行中の接触圧維持に関する力学的分析を組み合わせる検証タスクとして位置づけた．また，第2章では，皮むき動作が不規則な曲面に対する軌道追従と接触力維持の統合を必要とするタスクであることを整理した．皮むき動作では，ピーラーの刃を食材表面に沿って移動させ，皮を連続的に除去する．しかし，食材の形状には個体差があり，表面には曲率の変化や凹凸が存在するため，ピーラーが追従すべき軌道は食材ごとに異なる．さらに，刃先が表面から離れると皮むきが行えず，過度に押し付けると食材を傷つける可能性がある．したがって，本章では，皮むき動作を，食材表面形状に応じた刃先軌道の生成と，実行時における接触状態の維持を同時に解く課題として扱う．

本章では，食材表面の3D点群モデルに基づく基準軌道生成と，力覚フィードバックによる接触補正を組み合わせることで，食材形状に適応した皮むき動作を生成する．具体的には，第3章で取得した食材表面の3D点群モデルを入力とし，食材の断面輪郭からピーラーが追従する基準軌道を抽出する．次に，表面法線の変化に基づいて基準軌道を複数のセグメントに分割し，各セグメントに沿ってピーラーを移動させる．実行時には，力覚センサから得られる接触力を用いて刃先位置を補正し，点群モデルと実物表面の間に生じる誤差を吸収する．さらに，一回皮むき動作の後には，非皮領域と皮領域の境界を検出し，次に剥く領域がピーラー側に向くように食材を回転させる．



本節では，本章全体の前提となる問題設定，作業手順，および使用する点群情報を整理する．
まず，双腕ロボットによる皮むき作業を，ピーラー側のアームと食材固定側のアームの役割分担として定式化する．
次に，皮むき作業全体を，一回皮むき動作と食材の回転動作の反復として説明する．
最後に，本章で利用する食材表面点群，表面法線，RGB情報，作業台平面，および回転軸の役割を整理する．
なお，点群の取得，座標変換，多視点点群の統合，平滑化，および法線推定などの共通処理は第3章で述べたため，本章ではそれらの処理後の点群モデルを入力として扱う．
また，本章では皮がむかれていない領域を「皮領域」，ピーラーによって皮が除去された領域を「非皮領域」と定義する．


\subsection{問題設定}
\label{sec: peeling-problem-formulation}

皮むき動作の基本構成は，ピーラーを操作するアームと食材を支えるアームの役割分担にある．
本章では，片方のロボットアームがピーラーを固定し，もう片方のロボットアームが食材を支える双腕構成を用いる．
ピーラーを固定するアームをピーラーアーム，食材を支え回転させるアームを食材固定アームと呼ぶ．
ピーラーアームは，ピーラーの刃先を食材表面に接触させ，皮むき方向に沿って移動させることで，一部の皮を除去する．
食材固定アームは，一回の皮むき動作が終了した後，皮領域がピーラー側に向くように食材を回転させる．


本章で対象とする食材は，串に固定され，既知の回転軸まわりに回転できるものとして扱う．
対象とする食材の表面形状を \(S\) とし，食材の回転軸を \(\boldsymbol{v}_{\mathrm{rot}}\) とする．
食材固定アームは \(\boldsymbol{v}_{\mathrm{rot}}\) 周りに食材を回転させる役割を担う．
一方，ピーラーアームは，現在の食材姿勢に対して生成された基準軌道に沿って刃先を移動させる．
このように，皮むき作業は食材の回転とピーラーの移動に分けて扱うことができる．


皮むき動作を成立させる課題は，食材表面に沿う刃先軌道と，適切な接触状態を同時に実現することである．
食材表面は不規則な曲面であるため，軌道は一定ではない．
また，点群モデルから得られる表面位置には，RGB-Dカメラの計測誤差，多視点統合時のずれ，平滑化処理による形状変化が含まれる．
そのため，点群上で生成した軌道を位置制御のみで再生すると，ピーラーの刃が食材表面から離れて皮むきが行えない場合や，逆に過度に押し付けられる場合がある．
したがって，基準軌道の生成に加えて，実行中の接触力を用いた補正が必要となる．


本章で解く問題は，食材表面の3D点群モデルから，基準軌道，接触補正，および食材回転角を順に導出する問題として整理できる．
まず，食材表面点群から基準軌道を生成する．
次に，表面法線の変化に基づいて基準軌道を複数の軌道区間に分割し，実機で実行しやすい形に変換する．
さらに，ピーラーと食材表面との接触を維持するために，力覚フィードバックを用いて刃先位置を補正する．
一回皮むき動作の後には，非皮領域と皮領域の境界を検出し，その境界情報から次回の食材回転角を決定する．
最後に，皮むき動作と食材回転を反復し，終了条件を満たすまで皮むき範囲を拡大する．


この問題設定は，人間がピーラーを用いて食材を剥く基本的な手順を，双腕ロボット上で再構成するものである．
人間は，片方の手でピーラーを持ち，もう片方の手で食材を保持し，ピーラーを表面に沿って一方向に滑らせる．
その後，食材を少し回転させ，既に剥いた領域の隣接領域を同様に剥く．
本章の方法は，この手順を，一回皮むき動作と食材の回転操作の反復として表現し，それぞれに必要な幾何情報と力覚情報を導入するものである．


\subsection{皮むき動作の全体手順}

皮むき作業全体は，一回皮むき動作と食材の回転操作を終了条件まで繰り返す反復処理として構成される．
ここで，一回皮むき動作とは，ピーラーを食材表面に接触させ，あらかじめ生成した基準軌道に沿って移動させることで，食材表面の一部の皮を剥く処理である．
食材の回転操作とは，一回皮むき動作の後に，次に剥くべき皮領域がピーラー側へ向くように，食材を回転軸 \(\boldsymbol{v}_{\mathrm{rot}}\) 周りに回転させる処理である．
この二つの処理を交互に実行することで，局所的な皮むき動作を食材表面全体へ展開する．


各反復の前半では，現在の食材姿勢に対する基準軌道を食材表面点群から生成する．
まず，第3章で述べた手順により取得した食材表面のRGB点群 \(P_{\mathrm{all}}\) を入力する．
次に，食材の回転軸 \(\boldsymbol{v}_{\mathrm{rot}}\) を含む切断平面を設定し，食材表面点群から皮むき方向に対応する断面を抽出する．
この断面の輪郭は，ピーラーの刃先が追従する候補軌道となる．
ただし，断面輪郭のうち食材の下側や串の近傍は，作業台，エンドエフェクタ，またはロボットアームとの干渉が生じやすい．
そのため，到達困難な部分を除去し，到達可能な基準軌道 \(L\) を得る．


各反復の実行段階では，基準軌道を複数のセグメントに分割し，セグメントごとにピーラーを移動させる．
基準軌道 \(L\) は食材表面の曲率を反映した曲線であるが，実機で安定に実行するためには，曲線全体を一つの連続軌道として扱うよりも，短い軌道区間に分割した方が扱いやすい．
そこで，\(L\) 上の表面法線の変化を用いて，基準軌道を複数のセグメント \(s_j\) に分割する．
ピーラーアームは，各セグメントに沿って刃先を移動させ，セグメント内の平均法線に基づいてピーラー姿勢を設定する．


各反復の皮むき実行では，幾何的に生成された基準軌道に力覚フィードバックを重ねることで，接触力を維持する．
基準軌道は，食材表面の概略形状に基づく刃先の移動方向を与える．
しかし，実際の皮むきでは，ピーラーの刃が食材表面に適切に接触し続ける必要がある．
そこで，実行中には力覚センサの測定値を用いて接触状態を監視し，目標接触力との誤差に基づいて刃先位置を補正する．
これにより，点群計測誤差や表面凹凸による位置ずれを吸収しながら，ピーラーが皮に食い込む状態を維持する．


各反復の後半では，非皮領域と皮領域の境界を検出し，次回の食材回転角を計算する．
一回皮むき動作が終了した後，皮むき後の食材表面を観測し，RGB情報を用いて非皮領域を抽出する．
さらに，非皮領域と皮領域の境界を \(L_p\) として検出する．
この境界近傍の法線情報を平均することで平均法線 \(\boldsymbol{n}_{\mathrm{avg}}\) を求め，基準ベクトルとの関係から次回の食材回転角 \(\alpha\) を計算する．
食材固定アームは，この回転角に従って食材を \(\boldsymbol{v}_{\mathrm{rot}}\) 周りに回転させる．


本章の皮むき手順は，軌道生成，皮むき実行，境界抽出，食材回転，終了判定を閉ループ的に接続する点に特徴がある．
食材姿勢が更新されるたびに，現在観測されている表面に対して新たな基準軌道を生成し，一回皮むき動作を実行する．
その後，非皮領域の割合に基づいて作業の進行状況を判定し，終了条件を満たさない場合には次の反復へ進む．
このように，点群モデルに基づく幾何処理，実行中の力覚補正，および皮むき後の視覚認識を組み合わせることで，単一の皮むき動作を食材全体の皮むき作業へ拡張する．


\subsection{使用する点群情報}

本章で使用する点群情報は，皮むき動作の軌道生成，接触補正，回転角計算，および終了判定を結びつける基礎情報である．
点群の取得や前処理そのものは第3章で扱っているため，本章では取得済みの点群モデルから，皮むき動作に必要な情報を抽出して利用する．
具体的には，食材表面点群，表面法線，RGB情報，作業台平面，および食材の回転軸を用いる．
これらの情報は異なる処理で利用されるが，最終的にはピーラーの動作生成と食材固定アームの回転操作を接続する役割を担う．


食材表面点群 \(P_{\mathrm{all}}\) は，皮むき動作生成における最も基本的な入力である．
\(P_{\mathrm{all}}\) は，食材全体の外表面を表すRGB点群であり，食材形状と色情報を同時に含む．
本章では，\(P_{\mathrm{all}}\) から断面輪郭を抽出し，基準軌道を生成する．
また，一回皮むき動作後には，この点群に含まれる色情報を利用して非皮領域を検出する．
さらに，非皮領域の点群と対象表面点群の点数比を用いることで，皮むき作業の終了判定にも利用する．


表面法線は，食材表面の局所的な向きと曲率変化を表す幾何情報である．
点群中の点を \(\boldsymbol{p}_i\)，その点に対応する表面法線を \(\boldsymbol{n}_i\) とする．
表面法線は，基準軌道上の法線変化を用いた軌道分割に利用される．
また，各セグメントにおけるピーラー姿勢を決定する際にも，セグメント内の平均法線を用いる．
さらに，一回皮むき動作後に検出される非皮領域と皮領域の境界 \(L_p\) の近傍法線を平均することで，次回の食材回転角を決定するための情報を得る．


RGB情報は，非皮領域と皮領域を区別するための観測情報である．
皮むき動作では，食材の皮が残っている皮領域と，ピーラーにより皮が除去された非皮領域との間に色差が生じる場合がある．
本章では，この色差を利用して非皮領域を抽出する．
抽出された非皮領域は，次回の回転角を計算するための境界検出に用いられるだけでなく，作業が十分に進んだかを判定するための終了条件にも用いられる．
ただし，食材の種類によっては皮領域と非皮領域の色差が小さい場合があるため，RGB情報の利用には限界がある．


作業台平面と食材の回転軸は，点群上の幾何処理をロボット動作へ接続するための基準である．
作業台平面を \(P_T\)，その法線を \(\boldsymbol{n}_{\mathrm{ct}}\) とする．
また，食材の回転軸を \(\boldsymbol{v}_{\mathrm{rot}}\) とする．
作業台平面は，食材の下側やロボットが到達しにくい領域を判定するための基準となる．
一方，回転軸 \(\boldsymbol{v}_{\mathrm{rot}}\) は，断面抽出のための基準平面を定義する際に用いられるとともに，食材固定アームが食材を回転させる際の軸として用いられる．
したがって，\(P_T\)，\(\boldsymbol{n}_{\mathrm{ct}}\)，および \(\boldsymbol{v}_{\mathrm{rot}}\) は，基準軌道の生成と食材回転操作を結び付けるための共通基準として機能する．





% 第二章
\section{幾何情報による皮むき動作の分析}

本節では，食材表面の3D点群モデルから，皮むき動作に必要な幾何情報を抽出する方法について述べる．
前節で述べたように，本章では第3章で取得した食材表面点群を入力とし，点群の取得，座標変換，多視点点群の統合，平滑化，および法線推定などの共通処理は繰り返して述べない．
ここでは，処理後の食材表面点群 \(P_{\mathrm{all}}\) を用いて，ピーラーの刃先が倣う基準軌道 \(L\)，基準軌道を実行しやすく分割したセグメント \(s_j\)，および次回の食材回転角 \(\alpha\) を求める処理に焦点を当てる．


皮むき動作における幾何情報の役割は，食材の個体差をロボットが実行可能な動作パラメータへ変換することである．
同じ種類の食材であっても，形状，大きさ，表面曲率，および凹凸は個体ごとに異なる．
そのため，事前に固定した軌道だけを用いると，ピーラーの刃が食材表面から離れて皮むきが行えない場合や，反対に過度に押し付けられる場合がある．
そこで本章では，食材表面点群から断面輪郭を抽出し，その輪郭に基づいて一回皮むき動作の基準軌道を生成する．
さらに，基準軌道上の法線変化を用いて軌道を複数のセグメントに分割し，後続の動作計画で扱いやすい形に変換する．


また，皮むき作業全体を進めるためには，一回皮むき動作の後に，次に剥くべき皮領域をピーラー側へ向ける必要がある．
このため，本節では，皮むき後に得られるRGB点群から非皮領域と皮領域の境界を抽出し，その境界近傍の平均法線に基づいて食材の回転角を計算する方法についても述べる．
すなわち，本節で扱う幾何情報は，一回皮むき動作のための基準軌道生成と，次の一回皮むき動作へ移るための食材姿勢更新の両方を支える情報である．


\subsection{軌道の抽出}

一回皮むき動作では，ピーラーの刃先を食材表面に沿って移動させる必要がある．
この刃先の移動経路を，ここでは基準軌道と呼ぶ．
基準軌道は，食材表面の形状に応じて生成される幾何的な軌道であり，後続の力覚フィードバック制御によって実行時の接触誤差が補正される．
したがって，本項の目的は，点群モデルからピーラーが倣うべき表面上の概略的な移動方向を求めることである．


本研究では，一回皮むき動作でできるだけ広い領域を剥くため，ピーラーの刃先が食材の回転軸に沿う方向へ移動するように基準軌道を設定する．
食材の回転軸を \(\boldsymbol{v}_{\mathrm{rot}}\) とし，回転軸上の一点を \(\boldsymbol{p}_{\mathrm{rot}}\) とする．
また，作業台平面を \(P_T\)，その法線を \(\boldsymbol{n}_{\mathrm{ct}}\) とする．
食材表面から皮むき方向に対応する断面を得るため，まず，\(\boldsymbol{v}_{\mathrm{rot}}\) を含み，作業台平面 \(P_T\) に垂直な切断平面 \(P_{c1}\) を定義する．
この切断平面の法線 \(\boldsymbol{n}_{c1}\) は，次式で表される．


\[
\boldsymbol{n}_{c1}
=
\frac{
\boldsymbol{v}_{\mathrm{rot}} \times \boldsymbol{n}_{\mathrm{ct}}
}{
\left\|
\boldsymbol{v}_{\mathrm{rot}} \times \boldsymbol{n}_{\mathrm{ct}}
\right\|
}
\]

このとき，切断平面 \(P_{c1}\) は次式で表される．


\[
P_{c1}
=
\left\{
\boldsymbol{p}
\mid
(\boldsymbol{p}-\boldsymbol{p}_{\mathrm{rot}})
\cdot
\boldsymbol{n}_{c1}
=0
\right\}
\]

食材形状 \(S\) をこの切断平面で切断した断面を \(A_c\)，その断面の外周に対応する輪郭を \(L_c\) とする．
ただし，実際の処理では連続的な食材形状 \(S\) を直接扱うのではなく，点群 \(P_{\mathrm{all}}\) を用いる．
そこで，点群中の各点 \(\boldsymbol{p}_i\) と切断平面 \(P_{c1}\) の距離に基づき，切断平面近傍の点を抽出する．
抽出幅を \(d_1\) とすると，断面輪郭を構成する候補点群 \(P_c\) は次式で表される．


\[
P_c
=
\left\{
\boldsymbol{p}_i \in P_{\mathrm{all}}
\mid
\left|
(\boldsymbol{p}_i-\boldsymbol{p}_{\mathrm{rot}})
\cdot
\boldsymbol{n}_{c1}
\right|
\leq d_1
\right\}
\]

ここで，\(d_1\) は点群密度に応じて設定する閾値である．
この処理により，食材表面と切断平面の交線に対応する点群が得られる．
本章では，この点群を断面輪郭 \(L_c\) の点群表現として扱う．


しかし，断面輪郭 \(L_c\) の全体をそのまま基準軌道として用いることはできない．
食材の下側部分は作業台に近く，ピーラー，エンドエフェクタ，またはロボットアームが作業台と干渉する可能性がある．
また，食材の裏側に相当する部分では，ピーラーの刃を安定に接触させることが困難である．
そこで，到達困難な部分を除去するため，第2の切断平面 \(P_{c2}\) を導入する．


切断平面 \(P_{c2}\) は，食材の回転軸 \(\boldsymbol{v}_{\mathrm{rot}}\) を含み，作業台平面 \(P_T\) と平行な平面として定義する．
まず，\(\boldsymbol{v}_{\mathrm{rot}}\) と \(\boldsymbol{n}_{\mathrm{ct}}\) の両方に直交する補助ベクトル \(\boldsymbol{n}_p\) を次式で定義する．


\[
\boldsymbol{n}_p
=
\frac{
\boldsymbol{n}_{\mathrm{ct}} \times \boldsymbol{v}_{\mathrm{rot}}
}{
\left\|
\boldsymbol{n}_{\mathrm{ct}} \times \boldsymbol{v}_{\mathrm{rot}}
\right\|
}
\]

この補助ベクトルを用いて，切断平面 \(P_{c2}\) の法線 \(\boldsymbol{n}_{c2}\) を次式で定義する．


\[
\boldsymbol{n}_{c2}
=
\frac{
\boldsymbol{v}_{\mathrm{rot}} \times \boldsymbol{n}_p
}{
\left\|
\boldsymbol{v}_{\mathrm{rot}} \times \boldsymbol{n}_p
\right\|
}
\]

なお，\(\boldsymbol{n}_{c2}\) の向きは，作業台平面の上側を正とするように，\(\boldsymbol{n}_{c2}\cdot\boldsymbol{n}_{\mathrm{ct}} \geq 0\) を満たす向きに選択する．
これにより，切断平面 \(P_{c2}\) は，食材の回転軸を通り，作業台平面と平行な到達可能性判定用の平面となる．


切断平面 \(P_{c2}\) は次式で表される．


\[
P_{c2}
=
\left\{
\boldsymbol{p}
\mid
(\boldsymbol{p}-\boldsymbol{p}_{\mathrm{rot}})
\cdot
\boldsymbol{n}_{c2}
=0
\right\}
\]

この平面を用いて，断面輪郭 \(L_c\) のうち下側に存在する点を除去する．
本章では，次式を満たす点を到達可能な側の点として残す．


\[
(\boldsymbol{p}_i-\boldsymbol{p}_{\mathrm{rot}})
\cdot
\boldsymbol{n}_{c2}
\geq 0
\]

したがって，到達可能な基準軌道を構成する点群 \(P_L\) は次式で得られる．


\[
P_L
=
\left\{
\boldsymbol{p}_i \in P_c
\mid
(\boldsymbol{p}_i-\boldsymbol{p}_{\mathrm{rot}})
\cdot
\boldsymbol{n}_{c2}
\geq 0
\right\}
\]

この点群 \(P_L\) により表される曲線を，最終的な基準軌道 \(L\) とする．
この処理により，食材表面の断面輪郭のうち，ロボットが到達しやすい上側部分のみが基準軌道として残る．
得られた基準軌道 \(L\) は，食材の個体差を反映した一回皮むき動作の幾何的な基準であり，後続の軌道分割と動作計画に用いられる．


\subsection{法線による軌道分割}

前項で得られた基準軌道 \(L\) は，食材表面の曲率に従う曲線である．
一方，実機でピーラーを安定に移動させるためには，この曲線をそのまま一つの連続軌道として扱うよりも，複数の短い軌道区間として扱う方が実装しやすい．
また，ピーラーの刃は食材表面に対して受動的に向きを変えられるため，局所的な範囲では曲線軌道を短い直線セグメントとして近似できる．
そこで，本研究では，基準軌道 \(L\) 上の表面法線の変化に基づいて軌道を分割する．


まず，基準軌道 \(L\) を構成する点群を \(P_L\) とする．
点群として得られた \(P_L\) は，必ずしもピーラーが移動すべき順序で保存されているとは限らない．
そのため，軌道分割の前処理として，点群を皮むき方向に沿った順に並べる．
本章では，回転軸 \(\boldsymbol{v}_{\mathrm{rot}}\) に沿う成分を用いて，各点 \(\boldsymbol{p}_i\) の軌道方向座標 \(\lambda_i\) を次式で定義する．


\[
\lambda_i
=
(\boldsymbol{p}_i-\boldsymbol{p}_{\mathrm{rot}})
\cdot
\frac{
\boldsymbol{v}_{\mathrm{rot}}
}{
\left\|
\boldsymbol{v}_{\mathrm{rot}}
\right\|
}
\]

この \(\lambda_i\) の大小関係に基づいて \(P_L\) を並べ替えることで，ピーラーが始点から終点へ移動する順序を得る．
並べ替え後の点列を次式で表す．


\[
P_L
=
\left\{
\boldsymbol{p}_1,
\boldsymbol{p}_2,
\ldots,
\boldsymbol{p}_N
\right\}
\]

また，各点 \(\boldsymbol{p}_i\) に対応する表面法線を \(\boldsymbol{n}_i\) とし，法線集合を次式で表す．


\[
N_L
=
\left\{
\boldsymbol{n}_1,
\boldsymbol{n}_2,
\ldots,
\boldsymbol{n}_N
\right\}
\]

軌道分割では，現在のセグメントの開始点を \(\boldsymbol{p}_{\mathrm{st}}\) とし，その点に対応する法線を基準法線 \(\boldsymbol{n}_{\mathrm{ref}}\) とする．
点列に沿って各点を順に調べ，点 \(\boldsymbol{p}_i\) の法線 \(\boldsymbol{n}_i\) と基準法線 \(\boldsymbol{n}_{\mathrm{ref}}\) のなす角 \(\phi_i\) を計算する．


\[
\phi_i
=
\cos^{-1}
\left(
\frac{
\boldsymbol{n}_i \cdot \boldsymbol{n}_{\mathrm{ref}}
}{
\left\|
\boldsymbol{n}_i
\right\|
\left\|
\boldsymbol{n}_{\mathrm{ref}}
\right\|
}
\right)
\]

この角度 \(\phi_i\) が閾値角度 \(\theta_{\mathrm{th}}\) を超える場合，現在のセグメント内では表面法線の変化が大きくなったと判断し，その直前までの点列を一つのセグメントとして確定する．
すなわち，次式を満たす点を検出した位置を，セグメント境界として扱う．


\[
\phi_i > \theta_{\mathrm{th}}
\]

具体的には，現在のセグメント開始点のインデックスを \(k\) とし，点列を \(i=k+1,k+2,\ldots\) の順に調べる．
はじめに \(\boldsymbol{p}_{\mathrm{st}}=\boldsymbol{p}_k\)，\(\boldsymbol{n}_{\mathrm{ref}}=\boldsymbol{n}_k\) とする．
ある点 \(\boldsymbol{p}_i\) において \(\phi_i > \theta_{\mathrm{th}}\) を満たした場合，\(\boldsymbol{p}_k\) から \(\boldsymbol{p}_{i-1}\) までを一つのセグメント \(s_j\) として定義する．
その後，\(\boldsymbol{p}_i\) を次のセグメントの開始点とし，\(\boldsymbol{n}_i\) を新しい基準法線として同じ処理を繰り返す．


この処理により，基準軌道 \(L\) は複数のセグメントに分割される．


\[
L
=
\left\{
s_1,
s_2,
\ldots,
s_m
\right\}
\]

各セグメント \(s_j\) は，局所的には表面法線の変化が小さい軌道区間である．
そのため，後続の動作計画では，各セグメントを短い直線軌道として扱い，セグメント内の平均法線に基づいてピーラーの姿勢を決定する．
セグメント \(s_j\) に含まれる点集合を \(P_{s_j}\) とすると，その平均法線 \(\bar{\boldsymbol{n}}_j\) は次式で求められる．


\[
\bar{\boldsymbol{n}}_j
=
\frac{
\sum_{\boldsymbol{p}_i \in P_{s_j}}
\boldsymbol{n}_i
}{
\left\|
\sum_{\boldsymbol{p}_i \in P_{s_j}}
\boldsymbol{n}_i
\right\|
}
\]

この平均法線は，ピーラーの刃を食材表面へ向けるための代表的な表面方向として用いる．
点ごとの法線に逐次的に追従すると，点群ノイズや局所的な凹凸により，ピーラー姿勢が不安定になる可能性がある．
これに対して，セグメント単位で平均法線を用いることで，局所的なノイズの影響を抑えつつ，食材表面の大まかな曲率変化に応じた動作を生成できる．


以上のように，法線による軌道分割は，点群モデルから得られた曲線状の基準軌道を，ロボットが安定に実行しやすい動作単位へ変換する処理である．
表面法線方向の変化が大きい食材ではセグメント数が増え，法線変化が小さい食材では少数のセグメントで軌道を表現できる．
したがって，軌道分割は，食材形状の違いを動作計画の複雑さへ反映する役割を持つ．


\subsection{皮むき済み境界に基づく回転角の計算}

一回皮むき動作の後，次に剥くべき皮領域をピーラー側へ向けるためには，食材を適切な角度だけ回転させる必要がある．
この回転角は，直前の一回皮むき動作によって生成された非皮領域と，まだ皮が残っている皮領域との境界から求める．
境界は，次に剥くべき領域に隣接する幾何的な手がかりであるため，この境界の表面方向を用いることで，次回の食材姿勢を決定できる．


まず，一回皮むき動作の後に，皮むき後の食材表面を観測し，RGB情報を含む点群 \(P_{\mathrm{rgb}}\) を取得する．
点 \(\boldsymbol{p}_i\) の色情報を \(C(\boldsymbol{p}_i)\) とし，あらかじめ定義した皮領域の色範囲を \(C_{\mathrm{skin}}\) とする．
皮領域の色範囲に含まれない点を非皮領域の候補として抽出すると，非皮領域点群 \(P_{\mathrm{non}}\) は次式で表される．


\[
P_{\mathrm{non}}
=
\left\{
\boldsymbol{p}_i \in P_{\mathrm{rgb}}
\mid
C(\boldsymbol{p}_i) \notin C_{\mathrm{skin}}
\right\}
\]

この処理により，ピーラーによって皮が除去された領域を点群として得る．
ただし，色のみを用いた領域抽出は，照明条件，食材表面の色むら，および皮領域と非皮領域の色差に影響される．
そのため，本章では非皮領域全体を直接回転角へ変換するのではなく，非皮領域と皮領域の境界に注目する．


次に，非皮領域点群 \(P_{\mathrm{non}}\) から，回転方向側の端に相当する点を求める．
以下では，皮むき済み領域が次に回転する方向を \(+x\) 方向として表す．
このとき，非皮領域の端点の \(x\) 座標 \(X_{\mathrm{pe}}\) を次式で定義する．


\[
X_{\mathrm{pe}}
=
\max_{\boldsymbol{p}_i \in P_{\mathrm{non}}}
x_i
\]

ここで，\(x_i\) は点 \(\boldsymbol{p}_i\) の \(x\) 座標である．
また，この \(X_{\mathrm{pe}}\) を持つ点を，非皮領域の端点 \(\boldsymbol{p}_{\mathrm{pe}}\) とする．
実装上は，同じ \(x\) 座標を持つ点が複数存在する場合，観測側に近い点，または \(z\) 座標が大きい点を代表点として選択する．


次に，\(X_{\mathrm{pe}}\) を基準として，非皮領域と皮領域の境界候補点を抽出する．
境界抽出幅を \(d_2\) とすると，皮むき済み境界を構成する点群 \(P_{L_p}\) は次式で表される．


\[
P_{L_p}
=
\left\{
\boldsymbol{p}_i \in P_{\mathrm{rgb}}
\mid
X_{\mathrm{pe}}
\leq
x_i
\leq
X_{\mathrm{pe}}+d_2
\right\}
\]

この点群により表される境界を \(L_p\) とする．
\(d_2\) は，境界近傍の点を安定に取得するための幅であり，点群密度およびRGB情報のばらつきに応じて設定する．
この処理により，直前に剥いた非皮領域の端部に位置し，次に剥くべき皮領域に隣接する境界点群が得られる．


回転角は，境界 \(L_p\) 近傍の平均法線と基準ベクトルのなす角として計算する．
境界点群 \(P_{L_p}\) に含まれる各点の表面法線を \(\boldsymbol{n}_i\) とする．
境界近傍で用いる点数を \(N_{\mathrm{pe}}\) とすると，\(N_{\mathrm{pe}}\) は次式で表される．


\[
N_{\mathrm{pe}}
=
\left|
P_{L_p}
\right|
\]

境界近傍の平均法線 \(\boldsymbol{n}_{\mathrm{avg}}\) は，次式で求められる．


\[
\boldsymbol{n}_{\mathrm{avg}}
=
\frac{
\sum_{\boldsymbol{p}_i \in P_{L_p}}
\boldsymbol{n}_i
}{
\left\|
\sum_{\boldsymbol{p}_i \in P_{L_p}}
\boldsymbol{n}_i
\right\|
}
\]

この平均法線は，非皮領域と皮領域の境界における代表的な表面方向を表す．
一方，次の一回皮むき動作を開始するために望ましい表面方向を，基準ベクトル \(\boldsymbol{v}_{\mathrm{ref}}\) とする．
\(\boldsymbol{v}_{\mathrm{ref}}\) は，ピーラー側から見た食材表面の基準的な向きであり，食材を回転させた後に境界近傍の表面法線が向くべき方向を表す．


したがって，食材の回転角 \(\alpha\) は，\(\boldsymbol{n}_{\mathrm{avg}}\) と \(\boldsymbol{v}_{\mathrm{ref}}\) のなす角として次式で求められる．


\[
\alpha
=
\cos^{-1}
\left(
\frac{
\boldsymbol{v}_{\mathrm{ref}}
\cdot
\boldsymbol{n}_{\mathrm{avg}}
}{
\left\|
\boldsymbol{v}_{\mathrm{ref}}
\right\|
\left\|
\boldsymbol{n}_{\mathrm{avg}}
\right\|
}
\right)
\]

ここで，食材の回転方向は，非皮領域に隣接する皮領域がピーラー側へ移動する向きとしてあらかじめ定める．
そのため，\(\alpha\) はその方向に沿った回転量として扱う．
食材固定アームは，この回転角 \(\alpha\) に基づいて，食材を回転軸 \(\boldsymbol{v}_{\mathrm{rot}}\) 周りに回転させる．
これにより，直前に剥いた領域の隣接領域がピーラー側へ移動し，次の一回皮むき動作を実行できる状態となる．


以上の処理により，皮むき後に得られるRGB情報と表面法線情報を用いて，次回の食材姿勢を決定できる．
RGB情報は非皮領域と皮領域の境界抽出に用いられ，法線情報はその境界の表面方向を表すために用いられる．
したがって，皮むき済み境界に基づく回転角計算は，一回皮むき動作の結果を次の食材回転操作へ接続する処理である．
ただし，皮領域と非皮領域の色差が小さい食材では，境界抽出が不安定になる場合がある．
この点は，後述する実験において，食材の種類ごとの結果として評価する．


本節で述べた処理により，食材表面点群から基準軌道 \(L\)，分割後のセグメント \(s_j\)，各セグメントの平均法線 \(\bar{\boldsymbol{n}}_j\)，皮むき済み境界 \(L_p\)，および食材回転角 \(\alpha\) が得られる．
次節では，これらの幾何情報を用いて，実際のピーラー動作と食材回転操作をどのように生成し，力覚フィードバックを組み合わせて実行するかについて述べる．


% 第三章
\section{皮むき動作計画・制御}

本節では，前節で得られた幾何情報を用いて，実際のロボット動作を生成し，実行中の接触状態を制御する方法について述べる．前節では，食材表面点群から到達可能な基準軌道 \(L\) を抽出し，さらに表面法線の変化に基づいて複数のセグメント \(s_j\) に分割した．これに対し，本節では，各セグメントをロボットが実行可能な手先軌道へ変換し，ピーラーと食材表面の接触力を維持しながら皮むき動作を行う．

皮むき動作では，幾何情報だけで動作を完全に決定することは困難である．その理由は，点群モデルから得られる表面位置が，RGB-Dカメラの計測誤差，点群統合時のずれ，および表面平滑化処理の影響を含むためである．また，食材表面には凹凸や硬さの局所差が存在し，ピーラーの刃が食材表面に接触した後の反力も一定ではない．したがって，本研究では，3D点群モデルに基づく基準軌道を大域的な動作計画として用い，実行中には力覚フィードバックにより刃先位置を補正する．これにより，食材形状に応じた軌道生成と，実接触に基づくロバストな実行を統合する．

本節で扱う処理は，一回皮むき動作の計画，接触力フィードバック制御，食材の回転操作，および反復実行と終了判定で構成される．一回皮むき動作では，分割済みの基準軌道を手先軌道へ変換し，ピーラーの接近，皮むき，離脱を行う．接触力フィードバック制御では，目標接触力と力覚センサの測定値の誤差に基づいて，ピーラーの押し付け方向の速度を補正する．食材の回転操作では，一回皮むき動作後に得られた非皮領域の境界情報に基づいて，次に剥くべき皮領域がピーラー側へ向くように食材を回転させる．最後に，非皮領域の割合に基づいて終了判定を行い，条件を満たさない場合には同じ処理を繰り返す．

\subsection{一回皮むき動作の計画}

一回皮むき動作では，基準軌道 \(L\) を分割して得られた各セグメント \(s_j\) に対して，ピーラーアームの手先軌道を生成する．ここで，\(j=1,\ldots,m\) とし，\(m\) は基準軌道から得られたセグメント数を表す．各セグメントは，食材表面上の短い軌道区間であり，ピーラーの刃先をその区間に沿って移動させることで，局所的な皮むき動作を実行する．

第 \(j\) セグメント \(s_j\) の始点を \(\boldsymbol{p}^{j}_{\mathrm{st}}\)，終点を \(\boldsymbol{p}^{j}_{\mathrm{ed}}\) とする．このとき，セグメントの皮むき方向を表す単位ベクトル \(\boldsymbol{t}_j\) は次式で定義する．

\[
\boldsymbol{t}_j
=
\frac{
\boldsymbol{p}^{j}_{\mathrm{ed}}
-
\boldsymbol{p}^{j}_{\mathrm{st}}
}{
\left\|
\boldsymbol{p}^{j}_{\mathrm{ed}}
-
\boldsymbol{p}^{j}_{\mathrm{st}}
\right\|
}
\]

また，セグメント \(s_j\) に含まれる点群の法線集合を \(\mathcal{N}_j\) とし，その平均法線を \(\bar{\boldsymbol{n}}_j\) とする．平均法線は次式で求める．

\[
\bar{\boldsymbol{n}}_j
=
\frac{
\sum_{\boldsymbol{n}_i \in \mathcal{N}_j}
\boldsymbol{n}_i
}{
\left\|
\sum_{\boldsymbol{n}_i \in \mathcal{N}_j}
\boldsymbol{n}_i
\right\|
}
\]

ピーラーの姿勢は，この平均法線 \(\bar{\boldsymbol{n}}_j\) と皮むき方向 \(\boldsymbol{t}_j\) に基づいて決定する．具体的には，ピーラーの刃が食材表面に接触し，刃先が \(\boldsymbol{t}_j\) 方向へ移動できるようにツール座標系を設定する．ピーラーの刃は食材表面に対して受動的に回転できるため，曲線全体を厳密な連続曲線として追従する必要はない．そこで，本研究では，基準軌道 \(L\) を複数の短い直線セグメントとして扱い，セグメントごとに平均法線を用いて姿勢を更新する．

各セグメントに対して，接近点と離脱点を設定する．接近点は，ピーラーが食材表面に接触する前に通過する点であり，セグメント始点から表面法線方向へ一定距離だけ離した位置とする．接近距離を \(d_{\mathrm{app}}\) とすると，接近点 \(\boldsymbol{p}^{j}_{\mathrm{app}}\) は次式で表される．

\[
\boldsymbol{p}^{j}_{\mathrm{app}}
=
\boldsymbol{p}^{j}_{\mathrm{st}}
+
d_{\mathrm{app}}\bar{\boldsymbol{n}}_j
\]

同様に，離脱点は，セグメント終点で皮むき動作を終えた後にピーラーを食材表面から離すための点である．離脱距離を \(d_{\mathrm{ret}}\) とすると，離脱点 \(\boldsymbol{p}^{j}_{\mathrm{ret}}\) は次式で表される．

\[
\boldsymbol{p}^{j}_{\mathrm{ret}}
=
\boldsymbol{p}^{j}_{\mathrm{ed}}
+
d_{\mathrm{ret}}\bar{\boldsymbol{n}}_j
\]

したがって，各セグメントにおける一回皮むき動作は，接近点への移動，接触開始，セグメントに沿った皮むき，および離脱点への移動から構成される．接近点から始点までは位置制御により移動し，ピーラーの刃を食材表面へ近づける．始点近傍で接触力が検出された後，ピーラーはセグメント方向 \(\boldsymbol{t}_j\) に沿って移動する．この移動中には，次項で述べる力覚フィードバックにより，表面法線方向の位置補正を行う．終点に到達した後，ピーラーを離脱点へ移動させ，次のセグメントへ移る．

このように，一回皮むき動作では，点群モデルから得られる軌道をそのまま再生するのではなく，ロボットが安定に実行できる接近，接触，皮むき，離脱の動作列へ変換する．これにより，複雑な曲面形状を持つ食材であっても，局所的には単純なセグメント追従問題として扱うことができる．

\subsection{接触力フィードバック制御}

3D点群モデルから生成した基準軌道は，食材表面上の幾何学的な軌道である．しかし，実際に皮を剥くためには，ピーラーの刃が食材表面に一定の力で接触し，皮に適切に食い込む必要がある．位置制御のみで基準軌道を追従した場合，点群計測誤差や食材表面の凹凸により，ピーラーが食材表面から離れて皮むきが行えない場合がある．逆に，表面内側へ過度に押し込まれると，食材本体を削りすぎる可能性がある．そこで，本研究では，基準軌道に沿った移動に力覚フィードバックを重ねることで，接触力を一定に保つ．

目標接触力を \(f_c\)，力覚センサにより測定された接触方向の力を \(f_s\) とする．このとき，接触力誤差 \(f_e\) を次式で定義する．

\[
f_e
=
f_c
-
f_s
\]

ここで，\(f_e>0\) の場合は測定された接触力が目標値より小さく，ピーラーを食材表面へ近づける必要があることを表す．一方，\(f_e<0\) の場合は接触力が目標値より大きく，ピーラーを食材表面から離す必要があることを表す．この誤差に基づき，接触方向の補正速度 \(v_z\) を次式で求める．

\[
v_z
=
a|f_e|
\left(
\frac{1}{1+\exp(-b f_e)}
-
\frac{1}{2}
\right)
\]

ここで，\(a\) は力誤差に対する速度の感度を表すパラメータであり，\(b\) は誤差がゼロ近傍にある場合の不感帯の幅を調整するパラメータである．この関数を用いることで，誤差の符号に応じて補正速度の向きが決まり，誤差の大きさに応じて補正速度が変化する．したがって，接触力が不足している場合にはピーラーを押し付ける方向へ移動させ，接触力が過大な場合にはピーラーを離す方向へ移動させることができる．

ただし，\(f_e\) がゼロに近い場合，式から得られる速度が非常に小さくなり，ロボットの分解能や摩擦の影響により実際の補正動作が生じにくい場合がある．そこで，実装では速度下限値 \(v_{\mathrm{lim}}\) を設定し，最終的な補正速度 \(v_{fz}\) を次式で決定する．

\[
v_{fz}
=
\begin{cases}
v_z, & v_z=0 \\
\mathrm{sign}(v_z)
\max
\left(
|v_z|,v_{\mathrm{lim}}
\right),
& v_z \neq 0
\end{cases}
\]

この補正速度を，セグメント方向への送り速度と合成することで，ピーラーアームの手先速度を決定する．セグメント方向の送り速度を \(v_t\) とし，接触補正方向を \(\boldsymbol{e}_z\) とすると，手先速度 \(\boldsymbol{v}_{\mathrm{cmd}}\) は次式のように表される．

\[
\boldsymbol{v}_{\mathrm{cmd}}
=
v_t \boldsymbol{t}_j
+
v_{fz}\boldsymbol{e}_z
\]

第1項は，皮むき方向に沿ってピーラーを移動させるための速度成分である．第2項は，接触力を目標値へ近づけるための補正速度成分である．このように，皮むき動作中の手先速度は，幾何情報から決まる接線方向の移動と，力覚情報から決まる法線方向の補正の和として与えられる．

本制御の役割は，3D点群モデルに基づく軌道生成の誤差を，実行時の接触情報によって吸収することである．基準軌道 \(L\) は，食材の大域的な形状に基づいてピーラーが進むべき方向を与える．一方，力覚フィードバックは，点群モデルと実表面の微小なずれ，および食材表面の局所的な凹凸に対して，刃先位置を逐次補正する．これにより，ピーラーが食材表面に接触し続ける状態を維持しながら，皮むき方向に沿って安定に移動できる．

\subsection{食材の回転操作}

一回皮むき動作が終了すると，ピーラーが通過した部分は非皮領域となる．次の皮むき動作では，この非皮領域の隣に位置する皮領域をピーラー側へ向ける必要がある．そのため，食材固定アームにより食材を回転軸 \(\boldsymbol{v}_{\mathrm{rot}}\) 周りに回転させる．本研究では，前節で述べた非皮領域の境界 \(L_p\) と平均法線 \(\boldsymbol{n}_{\mathrm{avg}}\) から求めた回転角 \(\alpha\) を用いて，食材の回転操作を行う．

食材の回転操作は，ピーラーアームによる皮むき動作とは分離して実行する．まず，ピーラーアームを食材表面から離脱させ，食材固定アームが安全に回転できる状態にする．次に，食材固定アームは，串を保持した状態で，回転軸 \(\boldsymbol{v}_{\mathrm{rot}}\) 周りに角度 \(\alpha\) だけ食材を回転させる．このとき，回転変換を \(R(\boldsymbol{v}_{\mathrm{rot}},\alpha)\) とすると，第 \(k\) 回目の反復における食材姿勢から，第 \(k+1\) 回目の食材姿勢への更新は，概念的に次式で表される．

\[
T^{k+1}_{\mathrm{food}}
=
R(\boldsymbol{v}_{\mathrm{rot}},\alpha)
T^{k}_{\mathrm{food}}
\]

ここで，\(T^{k}_{\mathrm{food}}\) は第 \(k\) 回目の反復における食材姿勢を表す．実際のロボット制御では，この回転を食材固定アームのツール座標系における回転動作として実行する．食材は串に固定されているため，串の軸を回転軸として用いることで，食材を大きく持ち替えることなく，次の皮むき姿勢へ移行できる．

回転角 \(\alpha\) は，非皮領域の境界を基準として決定されるため，固定角で食材を回す場合と比較して，既に剥いた領域と次に剥く領域の重なりを抑えることができる．特に，食材の半径や表面曲率が一様でない場合には，固定回転角を用いると，剥き残しや重複した皮むきが生じやすい．これに対し，境界 \(L_p\) の平均法線を用いることで，現在の非皮領域の位置に応じて回転量を決定できる．

一方で，非皮領域と皮領域の色差が小さい食材では，境界 \(L_p\) の抽出が不安定になる可能性がある．この場合，平均法線に基づく回転角の計算も不安定となるため，固定回転角 \(\alpha_{\mathrm{fix}}\) を代替的に用いる．固定回転角を用いる場合，食材形状に応じた最適な回転は行えないが，境界検出が困難な条件でも反復的な皮むき作業を継続できる．

食材の回転後には，現在の食材姿勢に対して再び点群を取得し，基準軌道を更新する．これは，食材の姿勢が変化すると，ピーラー側から見た表面形状と到達可能領域が変化するためである．したがって，本研究の皮むき作業は，一度生成した軌道を最後まで繰り返し利用するのではなく，各回転後の食材姿勢に応じて軌道を再生成する構成をとる．

\subsection{ループ実行と終了判定}

皮むき作業全体は，一回皮むき動作と食材の回転操作を反復することで実現される．各反復では，現在の食材姿勢に対して点群を取得し，基準軌道を生成し，セグメント列に分割し，力覚フィードバックを用いて皮むき動作を実行する．その後，非皮領域を抽出し，回転角を計算して食材を回転させる．この処理を終了条件が満たされるまで繰り返す．

終了判定には，現在観測されている側面における非皮領域の点数と，対象表面点数の比を用いる．非皮領域の点数を \(N_{\mathrm{pe}}\)，対象表面領域の点数を \(N_{\mathrm{side}}\) とすると，終了判定比率 \(r_s\) は次式で定義される．

\[
r_s
=
\frac{
N_{\mathrm{pe}}
}{
N_{\mathrm{side}}
}
\]

この値が閾値 \(r_{\mathrm{th}}\) 以上となった場合，現在の側面に対する皮むきが十分に進行したと判断し，皮むき作業を終了する．終了条件は次式で表される．

\[
r_s
\geq
r_{\mathrm{th}}
\]

一方，\(r_s<r_{\mathrm{th}}\) の場合には，まだ十分な範囲の皮が除去されていないと判断し，食材の回転操作を行った後，次の反復へ移る．この反復処理により，一回皮むき動作で得られる局所的な非皮領域を，食材表面の広い範囲へ拡張することができる．

反復実行の流れは以下のように整理できる．まず，食材表面点群 \(P_{\mathrm{all}}\) を取得し，断面輪郭に基づいて基準軌道 \(L\) を抽出する．次に，法線変化に基づいて \(L\) を複数のセグメント \(s_j\) に分割する．ピーラーアームは，セグメントごとに皮むき動作を実行し，その間，力覚フィードバックによって接触力を維持する．一回皮むき動作が終了した後，皮むき後のRGB点群から非皮領域を抽出し，境界 \(L_p\) を求める．さらに，境界近傍の平均法線から回転角 \(\alpha\) を計算し，食材固定アームが食材を回転させる．最後に，非皮領域の割合 \(r_s\) を計算し，終了条件を満たすかを判定する．

終了判定の閾値 \(r_{\mathrm{th}}\) は，剥き残しと作業時間の関係を調整するパラメータである．閾値を小さくすると，作業は早く終了するが，皮領域が残る可能性が高くなる．一方，閾値を大きくすると，より広い範囲を剥くことができるが，すでに非皮領域となった部分を再度なぞる可能性や，串近傍などロボットが到達しにくい領域で反復が増える可能性がある．したがって，終了判定は単なる作業停止条件ではなく，皮むき率，作業時間，およびロボットの到達可能性を調整する要素として位置づけられる．

以上のように，本節で述べた動作計画・制御は，3D点群モデルから得た幾何情報を用いて大域的な基準軌道を生成し，実行中には力覚情報により局所的な誤差を補正し，一回皮むき動作後には視覚情報により次の食材姿勢を決定する構成である．この構成により，形状や曲率が異なる食材に対しても，皮むき動作を反復的に生成できる．

% 第四章
\section{実験}

本節では，提案した皮むき動作生成手法を双腕ロボットシステムに実装し，複数種類の食材に対して実機実験を行った結果について述べる．実験の目的は，3D点群モデルに基づく幾何軌道生成，法線変化に基づく軌道分割，力覚フィードバックによる接触維持，および非皮領域境界に基づく食材回転が，形状や表面性状の異なる食材に対して有効に機能するかを確認することである．

本章で扱う皮むき動作は，単に一度だけピーラーを移動させる操作ではなく，軌道生成，皮むき実行，非皮領域認識，食材回転，終了判定を反復する作業である．したがって，実験では，個々の皮むき動作が成立するかだけでなく，反復処理によって食材表面の広い範囲を剥けるかを評価する．

\subsection{概要}

実験には，形状，曲率，および色特徴が異なる食材として，リンゴ，きゅうり，なす，じゃがいも，およびにんじんを用いた．リンゴについては，大，中，小の3種類の大きさを用意し，各大きさに対して3回ずつ実験を行った．その他の食材については，それぞれ1回ずつ実験を行った．

これらの食材を選定した理由は，提案手法の各構成要素の有効性と限界を確認するためである．リンゴは表面曲率が大きく，断面輪郭に基づく軌道生成が重要となる食材である．きゅうりとなすは表面曲率が比較的小さく，法線変化が緩やかであるため，軌道分割後のセグメント実行の安定性を確認する対象である．じゃがいもは表面がでこぼこしており，点群モデルから得た軌道と実際の接触状態のずれが生じやすい．したがって，力覚フィードバックによる接触補正の有効性を確認する対象である．にんじんは，皮領域と非皮領域の色差が小さく，非皮領域の認識が困難な食材である．そのため，RGB情報に基づく境界抽出の限界と，固定回転角を用いる代替方法の有効性を確認する対象である．

実験では，まず食材に串を刺し，食材固定アームのエンドエフェクタに取り付けた．ピーラーは，ピーラーアーム側のエンドエフェクタに固定した．その後，ピーラーアームに取り付けられたRGB-Dカメラを用いて食材表面のRGB点群を取得し，基準軌道を生成した．基準軌道を法線変化に基づいて分割した後，各セグメントに沿ってピーラーを移動させ，力覚フィードバックにより接触力を維持しながら皮むき動作を実行した．一回皮むき動作の後には，RGB点群を再取得し，非皮領域を抽出した．非皮領域の境界から回転角を計算し，食材固定アームにより食材を回転させた．この処理を終了条件が満たされるまで繰り返した．

実験で用いた主なパラメータは以下の通りである．接触力の目標値は \(f_c=5.0\,\mathrm{N}\) とした．力覚フィードバックの速度決定式におけるパラメータは \(a=10.0\)，\(b=0.05\) とした．また，補正速度の下限値は \(v_{\mathrm{lim}}=0.1\,\mathrm{mm/s}\) とした．軌道分割に用いる閾値角度は \(\theta_{\mathrm{th}}=30^{\circ}\) とした．終了判定の閾値は \(r_{\mathrm{th}}=0.9\) とした．にんじんについては，非皮領域の認識が困難であったため，境界に基づく回転角計算の代わりに，固定回転角 \(\alpha_{\mathrm{fix}}=20^{\circ}\) を用いた．

\subsection{評価指標}

皮むき結果は，非皮領域の割合により評価した．本来であれば，食材表面全体における非皮領域の面積を計算することが望ましい．しかし，食材全体の完全な表面積を高精度に求めることは容易ではなく，また皮むき後の表面を単一視点から完全に観測することも困難である．そこで，本実験では，面積の代わりにRGB点群中の点数を用いて非皮領域の割合を近似的に評価した．

評価用点群として，二方向から取得したRGB点群を用いた．一方の点群を \(P_A\)，他方の点群を \(P_B\) とする．それぞれの点数を \(N_A\)，\(N_B\) とし，色情報により抽出された非皮領域の点数を \(N_{fA}\)，\(N_{fB}\) とする．このとき，皮むき率 \(r\) を次式で定義する．

\[
r
=
\frac{
N_{fA}
+
N_{fB}
}{
N_A
+
N_B
}
\]

ここで，\(r\) が大きいほど，観測された食材表面に対して広い範囲の皮が除去されたことを表す．なお，にんじんについては，皮領域と非皮領域の色差が小さく，点群の色フィルタによる非皮領域抽出が不安定であった．そのため，にんじんの皮むき率は目測により評価した．

また，皮むき率に加えて，皮むき動作の安定性を定性的に評価した．具体的には，ピーラーと食材表面の接触が維持されたか，食材回転後に次の皮むき領域が既に剥いた領域の隣に配置されたか，剥き残しが発生した場合にその原因がロボット干渉，接触力不足，境界認識の失敗，または終了判定の閾値設定に由来するかを確認した．

\subsection{実験結果}

リンゴに対する実験結果を表\ref{tab:peeling_apple}に示す．リンゴでは，大，中，小の3種類の大きさについて，それぞれ3回の実験を行った．大きいリンゴでは，非皮領域の割合は \(78.01\%\)，\(79.83\%\)，\(82.69\%\) であり，平均は \(80.18\%\) であった．中程度のリンゴでは，\(88.85\%\)，\(84.25\%\)，\(84.11\%\) であり，平均は \(85.74\%\) であった．小さいリンゴでは，\(85.93\%\)，\(90.93\%\)，\(88.62\%\) であり，平均は \(88.49\%\) であった．

\begin{table}[tb]
    \centering
    \caption{リンゴの皮むき結果}
    \label{tab:peeling_apple}
    \begin{tabular}{c|ccc}
        \hline
        試行 & 大[\%] & 中[\%] & 小[\%] \\
        \hline
        1 & 78.01 & 88.85 & 85.93 \\
        2 & 79.83 & 84.25 & 90.93 \\
        3 & 82.69 & 84.11 & 88.62 \\
        \hline
        平均 & 80.18 & 85.74 & 88.49 \\
        \hline
    \end{tabular}
\end{table}

リンゴは表面曲率が大きく，一回皮むき動作中に表面法線方向が大きく変化する．そのため，固定された直線軌道ではなく，点群モデルから断面輪郭を抽出し，法線変化に基づいて軌道を分割することが有効であった．一方で，大きいリンゴでは，串に近い領域やロボットアームとの干渉が生じやすい領域に剥き残しが発生した．これは，幾何的には皮むき対象領域であっても，実機の到達可能性やエンドエフェクタとの干渉によりピーラーが十分に接近できない領域が存在するためである．

きゅうり，なす，じゃがいも，およびにんじんに対する結果を表\ref{tab:peeling_other_foods}に示す．きゅうりの非皮領域の割合は \(94.31\%\)，なすは \(90.23\%\)，じゃがいもは \(88.72\%\) であった．にんじんについては，非皮領域の認識が困難であったため目測による評価であるが，約 \(95.00\%\) の皮が除去された．

\begin{table}[tb]
    \centering
    \caption{各食材の皮むき結果}
    \label{tab:peeling_other_foods}
    \begin{tabular}{c|cccc}
        \hline
        食材 & きゅうり & なす & じゃがいも & にんじん \\
        \hline
        非皮領域の割合[\%] & 94.31 & 90.23 & 88.72 & 95.00 \\
        \hline
    \end{tabular}
\end{table}

きゅうりとなすでは，リンゴと比較して表面法線方向の変化が小さいため，軌道分割後の各セグメントを安定に実行できた．特に，きゅうりでは細長い形状に対してピーラーをほぼ一定方向に移動させることができ，高い皮むき率が得られた．なすでも \(90\%\) を超える非皮領域の割合が得られたが，一部に剥き残しが見られた．これは，終了判定の閾値を \(r_{\mathrm{th}}=0.9\) と設定していたため，非皮領域の割合が閾値をわずかに上回った時点で作業が終了したことが一因であると考えられる．

じゃがいもは表面がでこぼこしており，点群モデル上の基準軌道と実際の接触状態との間に差が生じやすい食材である．それにもかかわらず，\(88.72\%\) の非皮領域割合が得られた．この結果は，力覚フィードバックによりピーラーの刃先位置が補正され，表面の凹凸に対して接触状態を維持できたことを示している．ただし，凹凸の大きい部分では接触力の変動が生じやすく，皮の除去量が局所的に不均一になる場合があった．

にんじんでは，皮領域と非皮領域の色差が小さく，RGB情報に基づいて境界 \(L_p\) を安定に抽出することが困難であった．そのため，境界に基づく回転角計算ではなく，固定回転角を用いて食材を回転させた．この条件でも，目測では約 \(95\%\) の皮を除去できた．この結果は，非皮領域の認識が困難な食材に対しても，固定回転角を用いることで皮むき作業自体は継続可能であることを示している．一方で，定量評価の信頼性や回転角の最適性という点では，境界認識を用いる方法に比べて制約が残る．

以上の結果より，提案手法は，曲率の大きいリンゴ，曲率変化の小さいきゅうりとなす，表面がでこぼこしたじゃがいも，および色認識が困難なにんじんに対して，いずれも広い範囲の皮むきを実現できた．特に，すべての食材種別において，平均的にはおおむね \(80\%\) 以上の非皮領域割合が得られた．

\subsection{考察}

実験結果から，まず，幾何情報に基づく軌道生成の有効性が確認できる．リンゴのように表面曲率が大きい食材では，食材表面に沿った軌道を事前に生成しなければ，ピーラーの刃が表面から離れたり，過度に食い込んだりする可能性がある．本研究では，回転軸を含む断面輪郭から基準軌道を抽出し，さらに法線変化に基づいて軌道を複数のセグメントに分割した．その結果，曲率変化の大きい食材に対しても，ロボットが実行可能な形で皮むき軌道を生成できた．

次に，力覚フィードバック制御の有効性が確認できる．じゃがいものように表面が粗い食材では，点群モデルから得られる軌道と実際の食材表面との間にずれが生じやすい．位置制御のみを用いた場合，このずれによりピーラーが空振りする場合や，逆に食材へ過度に押し込まれる場合がある．これに対し，本研究では，目標接触力と測定値の誤差に基づいてピーラーの押し付け方向の速度を補正した．この構成により，表面凹凸を含む食材に対しても，皮むき中の接触状態を維持できた．

また，非皮領域境界に基づく回転角計算は，食材表面上の皮むき済み領域に応じて次の皮むき位置を決定するために有効である．固定回転角のみを用いる場合，食材の大きさや局所的な曲率によって，既に剥いた領域との重なりや剥き残しが生じやすい．一方，境界 \(L_p\) と平均法線 \(\boldsymbol{n}_{\mathrm{avg}}\) を用いることで，現在の非皮領域の位置に応じて食材を回転できる．このため，リンゴ，きゅうり，なす，じゃがいものように非皮領域を認識できる食材では，皮むき動作と食材回転を連続的に接続できた．

一方で，RGB情報に基づく境界抽出には限界がある．にんじんのように皮領域と非皮領域の色差が小さい食材では，色フィルタによる非皮領域抽出が不安定となる．この場合，境界に基づく回転角計算が困難であり，固定回転角を用いる必要があった．今後は，RGB情報だけでなく，表面反射，テクスチャ，切削後の形状変化，または力覚情報を組み合わせることで，非皮領域の認識を安定化する必要がある．

さらに，システム上の制約として，串に近い領域やエンドエフェクタ周辺の領域では，ピーラーアームや食材固定アームとの干渉により剥き残しが発生しやすい．この問題は，食材を串で固定して回転させる構成に由来するものであり，人間がピーラーを用いる場合にも，保持部近傍は剥きにくい領域となる．したがって，この領域を完全に剥くためには，食材の持ち替え，串の挿入位置の変更，または皮むき後に別の工程を追加する必要がある．

終了判定の閾値設定も結果に影響する．なすの結果に見られるように，非皮領域の割合が閾値を超えた時点で作業を終了すると，一部の皮領域が残る可能性がある．閾値を高く設定すれば剥き残しを減らせるが，作業時間が長くなり，すでに剥いた領域を再度なぞる可能性が高くなる．したがって，終了判定は，単純な割合だけでなく，未処理領域の位置，到達可能性，およびロボット干渉の有無を考慮して設計する必要がある．

最後に，処理時間についても課題が残る．本手法では，各反復において点群取得，法線推定，軌道抽出，軌道分割，非皮領域認識を行うため，計算に時間を要する．特に，法線推定や点群処理は処理時間に大きく影響する．家庭環境で実用的な皮むき作業を実現するためには，点群処理の高速化，軌道再生成の効率化，および必要な再計測回数の削減が重要である．

以上より，提案手法は，食材表面の3D点群モデルに基づく幾何的な軌道生成と，実行中の力覚フィードバック制御を組み合わせることで，形状や表面性状の異なる複数の食材に対して皮むき動作を実現できることを示した．一方で，色差が小さい食材における非皮領域認識，串近傍の剥き残し，終了判定の最適化，および処理時間の短縮が今後の課題である．

